%-*-latex-*-
NSA is designing a new steganographic system to hide information.
Here's how it works.

Suppose agent Bob wants to send a message \verb!mad! to Alice
and Eve the rogue agent is eavesdropping.
Bob appends the characters of the message, except the last,
to the message but in the reverse order to get \verb!madam!.
Note that this is a palindrome with odd length.
He then inserts a huge number of characters into the string to further
obfuscate the message.
Here's an example:
\[
  \texttt{hieve\underline{\textred{m}}yh\underline{\textred{a}}mburger\underline{\textred{d}}oesnott\underline{\textred{a}}stegoodwith\underline{\textred{m}}armalade}
\]
The problem is Bob might accidentally create many subsequences in
this obfuscated message.
So Bob and Alice has agreed that if there is more than one
possible subsequence, the longest one is the correct palindrome.
NSA has asked you to write a program  
to find the longest odd-length palindromic subsequence
of a given input string (hopefully there's only one).
Of course it's possible that there are more than one longest odd-length
palindrome subsequence.
In that case, your program will be helpful in warning Bob that his obfuscation is ambiguous.

(a) For a string of length $n$, how many subsequences are there altogether?
(Using brute force, you can then check each subsequence to see if the
length is odd and is it is a palindrome. Of course you know by now that brute force is a bad idea.)

(b) Define a suitable measure for optimization (a length function)
in order to solve this problem.
For instance if the input is \verb!x[0..n-1]!, maybe
you want the function $L(\verb!i!)$ to be the length of the longest palindromic
subsequence of \verb!x[0..i]!, or maybe
you want the length of the longest palindromic
subsequence of \verb!x[i..n-1]!
or maybe you want
the function is $L(\verb!i!, \verb!j!)$ to be the length of the longest
palindomic subsequence of \verb!x[i..j]!.
Or maybe you want something that is completely different (and more
crazy than what I can think of.)

(c)
Find a recursion on your function in (b).
Implement a top-down DP solution.

(d)
Find a bottom-up DP solution to our problem
and implement it in C\texttt{++}.
Here's a sample run:
\begin{Verbatim}[frame=single,commandchars=\\\{\}]
\userinput{whatacharacter}
("carac", (5, 7, 8, 9, 10))
\end{Verbatim}
There is only one input. (Bob wanted to say \verb!car!, so the palindrome is \verb!carac!.)
The output is the longest odd-length palindromic subsequence followed by the
indices of this subsequence.
If there are more than one longest palindromic subsequences,
print them on separate lines.
For instance in the above, there's another output:
\begin{Verbatim}[frame=single,commandchars=\\\{\}]
\userinput{whatacharacter}
("carac", (5, 7, 8, 9, 10))
("hatah", (1, 2, 3, 4, 6))
\end{Verbatim}
(Note that for this part, you should print the longest.
\verb!tcaract! is actually longer. The above test case
is only to fix the input-output format. Outputs shown above need not be correct.)

(e) What is the runtime and space complexity of your algorithm?
Write this in the documentation of your main.cpp:
\begin{Verbatim}[frame=single,fontsize=\small]
// T(n) = O(?)
// SPACE(n) = O(?)

... your program ...
\end{Verbatim}

(f) Finally using the above program
write another program that prints all odd-length palindromic subsequence
with length of at least 3.
Organized the output in ascending order by length.
For string with the same length, organize them by dictionary order.

\textsc{Note.}
Assume characters in string are in lowercase.
